
\documentclass{article}
\usepackage{colortbl}
\usepackage{makecell}
\usepackage{multirow}
\usepackage{supertabular}

\begin{document}

\newcounter{utterance}

\centering \large Interaction Transcript for game `cladder', experiment `full\_v1.5\_default', episode 9125 with refineonly{-}llama{-}3.1{-}8b{-}instruct.
\vspace{24pt}

{ \footnotesize  \setcounter{utterance}{1}
\setlength{\tabcolsep}{0pt}
\begin{supertabular}{c@{$\;$}|p{.15\linewidth}@{}p{.15\linewidth}p{.15\linewidth}p{.15\linewidth}p{.15\linewidth}p{.15\linewidth}}
    \# & \multicolumn{2}{c}{Player} && \multicolumn{2}{c}{Game Master} \\
    \hline

    \theutterance \stepcounter{utterance}  
    & & & \multicolumn{4}{p{0.6\linewidth}}{
        \cellcolor[rgb]{0.9,0.9,0.9}{
            \makecell[{{p{\linewidth}}}]{
                \texttt{\tiny{[P1$\langle$GM]}}
                \texttt{You are an expert in causal inference. The following question is not a typical commonsense query, but rather a meticulously designed question created by a professor specializing in causal inference, intended to assess the students' mastery of the course content.} \\
\texttt{Imagine a self{-}contained, hypothetical world with only the following conditions, and without any unmentioned factors or causal relationships: Kraz has a direct effect on pexu and rukz. Pexu has a direct effect on hwax. Hwax has a direct effect on rukz. Kraz is unobserved. For those who are not pexu and are not hwax, the probability of rukz is 19\%. For those who are not pexu and are hwax, the probability of rukz is 29\%. For those who are pexu and are not hwax, the probability of rukz is 20\%. For those who are pexu and are hwax, the probability of rukz is 31\%. For those who are not pexu, the probability of hwax is 16\%. For those who are pexu, the probability of hwax is 90\%. For those who are pexu, would it be less likely to see rukz if the individual was not pexu?} \\
\texttt{Start your answer with 'Yes' or 'No', followed by additional reasoning or evidence to support your explanation.} \\
            }
        }
    }
     \\ \\

    \theutterance \stepcounter{utterance}  
    & & \multicolumn{4}{p{0.6\linewidth}}{
        \cellcolor[rgb]{0.95,0.95,0.95}{
            \makecell[{{p{\linewidth}}}]{
                \texttt{\tiny{[P1$\leftrightarrow$P]}}
                \texttt{{[}DRAFT{]}} \\
\texttt{No.} \\
\\ 
\texttt{To address this question, we need to apply the principles of causal inference, specifically the concept of counterfactuals and potential outcomes. } \\
\\ 
\texttt{First, let's establish the notation: } \\
\\ 
\texttt{{-} P(Rukz|Pexu, Hwax) denotes the probability of Rukz given Pexu and Hwax.} \\
\texttt{{-} P(Rukz|Pexu) denotes the probability of Rukz given Pexu.} \\
\texttt{{-} P(Hwax|Pexu) denotes the probability of Hwax given Pexu.} \\
\\ 
\texttt{We are given the following probabilities:} \\
\\ 
\texttt{{-} P(Rukz|Not Pexu, Not Hwax) = 0.19} \\
\texttt{{-} P(Rukz|Not Pexu, Hwax) = 0.29} \\
\texttt{{-} P(Rukz|Pexu, Not Hwax) = 0.20} \\
\texttt{{-} P(Rukz|Pexu, Hwax) = 0.31} \\
\texttt{{-} P(Hwax|Not Pexu) = 0.16} \\
\texttt{{-} P(Hwax|Pexu) = 0.90} \\
\\ 
\texttt{We want to compare the probability of Rukz given Pexu and not Pexu. To do this, we need to calculate the counterfactual probability of Rukz given not Pexu, which is P(Rukz|Not Pexu).} \\
\\ 
\texttt{Using the law of total probability, we can calculate P(Rukz|Not Pexu) as follows:} \\
\\ 
\texttt{P(Rukz|Not Pexu) = P(Rukz|Not Pexu, Not Hwax) * P(Not Hwax|Not Pexu) + P(Rukz|Not Pexu, Hwax) * P(Hwax|Not Pexu)} \\
\texttt{= 0.19 * (1 {-} 0.16) + 0.29 * 0.16} \\
\texttt{= 0.19 * 0.84 + 0.29 * 0.16} \\
\texttt{= 0.1604 + 0.0464} \\
\texttt{= 0.2068} \\
\\ 
\texttt{Now, we can compare the probability of Rukz given Pexu and not Pexu. We are given that P(Rukz|Pexu) = 0.31 (which is the average of P(Rukz|Pexu, Not Hwax) and P(Rukz|Pexu, Hwax)).} \\
\\ 
\texttt{To determine if it is less likely to see Rukz if the individual was not Pexu, we need to compare P(Rukz|Not Pexu) and P(Rukz|Pexu). Since P(Rukz|Not Pexu) = 0.2068 and P(Rukz|Pexu) = 0.31, we can} \\
            }
        }
    }
    & \\ \\

    \theutterance \stepcounter{utterance}  
    & \multicolumn{4}{p{0.6\linewidth}}{
        \cellcolor[rgb]{0.561,0.737,0.561}{
            \makecell[{{p{\linewidth}}}]{
                \texttt{\tiny{[P1$\rangle$GM]}}
                \texttt{No.} \\
\\ 
\texttt{Since P(Rukz|Not Pexu) = 0.2068 and P(Rukz|Pexu) = 0.31, we can see that P(Rukz|Pexu) is greater than P(Rukz|Not Pexu). This suggests that it is less likely to see Rukz if the individual was not Pexu.} \\
            }
        }
    }
    & & \\ \\

    \theutterance \stepcounter{utterance}  
    & & & \multicolumn{2}{p{0.3\linewidth}}{
        \cellcolor[rgb]{0.9,0.9,0.9}{
            \makecell[{{p{\linewidth}}}]{
                \texttt{\tiny{[GM$|$GM]}}
                \texttt{No} \\
            }
        }
    }
    & & \\ \\

    \theutterance \stepcounter{utterance}  
    & & & \multicolumn{2}{p{0.3\linewidth}}{
        \cellcolor[rgb]{0.9,0.9,0.9}{
            \makecell[{{p{\linewidth}}}]{
                \texttt{\tiny{[GM$|$GM]}}
                \texttt{yes} \\
            }
        }
    }
    & & \\ \\

    \theutterance \stepcounter{utterance}  
    & & & \multicolumn{2}{p{0.3\linewidth}}{
        \cellcolor[rgb]{0.9,0.9,0.9}{
            \makecell[{{p{\linewidth}}}]{
                \texttt{\tiny{[GM$|$GM]}}
                \texttt{game\_result = LOSE} \\
            }
        }
    }
    & & \\ \\

\end{supertabular}
}

\end{document}
